% future/tm.tex
% mainfile: ../perfbook.tex
% SPDX-License-Identifier: CC-BY-SA-3.0

\section{事务内存}
\label{sec:future:Transactional Memory}
\epigraph{一切应该尽可能简单，但不能过于简单。}
	 {Albert Einstein，经由Louis Zukofsky}


在数据库之外使用事务的想法可以追溯到几十年
前~\cite{DBLomet1977SIGSOFT,Knight:1986:AMF:319838.319854,Herlihy93a}，
数据库事务和非数据库事务之间的关键区别在于，非数据库事务
去掉了定义数据库事务的``ACID''\footnote{
	原子性（Atomicity）、一致性（Consistency）、隔离性（Isolation）
	和持久性（Durability）。}
属性中的``D''。
在硬件中支持基于内存的事务，即``transactional memory''
（TM）的想法出现得更
晚~\cite{Herlihy93a}，但不幸的是，尽管也有一些类似的提案被提
出~\cite{JMStone93}，商用硬件对此类事务的支持并未立即出现。
不久之后，Shavit和Touitou提出了一种纯软件的
transactional memory (STM) 实现，它能够在商用硬件上运行，
尽管存在内存排序问题~\cite{Shavit95}。
这一提案沉寂了多年，也许是因为研究社区的注意力被\IXacrl{nbs}
（见\cref{sec:advsync:Non-Blocking Synchronization}）所吸引。

但到了世纪之交，TM开始受到更多
关注~\cite{Martinez01a,Rajwar01a}，到了这个十年的中期，
人们的兴趣只能用``白热化''来形
容~\cite{MauriceHerlihy2005-TM-manifesto.pldi,
DanGrossman2007TMGCAnalogy}，只有少数谨慎的
声音~\cite{Blundell2005DebunkTM,McKenney2007PLOSTM}。

TM背后的基本思想是原子地执行一段代码，使得其他线程（thread）看不到中间状态。
因此，TM的语义可以简单地通过将每个事务替换为一个可递归获取的
全局锁的获取和释放来实现，尽管性能和可扩展性会极其糟糕。
TM实现（无论是硬件还是软件）中固有的大部分复杂性在于
高效地检测并发事务何时可以安全地并行运行。
由于这种检测是动态进行的，冲突的事务可能会被中止或``回滚''，
而且在某些实现中，这种失败模式对程序员是可见的。

由于事务回滚的可能性随着事务大小的减小而降低，
TM对于小型基于内存的操作可能变得相当有吸引力，
例如用于栈、队列、哈希表（hash table）和搜索树的链表操作。
然而，目前对于大型事务，特别是包含I/O和进程创建等
非内存操作的事务，要证明其适用性则困难得多。
以下各节审视了``Transactional Memory
Everywhere''~\cite{PaulEMcKenney2009TMeverywhere}
这一宏大愿景当前面临的挑战。
\Cref{sec:future:Outside World}考察了与外部世界交互所面临的挑战，
\cref{sec:future:Process Modification}探讨了与进程修改原语的交互，
\cref{sec:future:Synchronization}探索了与其他同步原语的交互，
最后\cref{sec:future:Discussion}以一些讨论作为结尾。

\subsection{外部世界}
\label{sec:future:Outside World}

用\ppl{Donald}{Knuth}的智慧之言来说：

\begin{quote}
	许多计算机用户觉得输入和输出实际上不是``真正的编程''的一部分，
	它们仅仅是为了将信息输入和输出机器而（不幸地）
	必须做的事情。
\end{quote}

无论我们是否相信输入和输出是``真正的编程''，
事实是软件绝对必须与外部世界打交道。
因此，本节评析transactional memory的外部世界能力，
重点关注I/O操作、时间延迟和持久存储。

而这些与外部世界的交互正是TM无条件可组合性声明
破灭的基石。
是的，你可以组合事务，但前提是没有介入的TM不友好操作。
正如你可以组合基于锁的临界区（critical section），但前提是这样做不会引入死锁（deadlock）。

最终，在实际代码中TM是否真的比锁提供更好的可组合性，
这一点并不清楚。

\subsubsection{I/O操作}
\label{sec:future:I/O Operations}

可以在基于锁的临界区内执行I/O操作，
在持有\IX{hazard pointer}时，在sequence-locking读侧临界区内，
以及在userspace-RCU读侧临界区内执行I/O操作，
如果需要的话，甚至可以同时在所有这些情况下执行。
当你尝试从事务内部执行I/O操作时会发生什么？

根本问题在于事务可能会被回滚，例如由于冲突。
粗略地说，这要求给定事务内的所有操作都是可撤销的，
即执行两次操作与执行一次具有相同的效果。
不幸的是，I/O通常是典型的不可撤销操作，
这使得将通用I/O操作包含在事务中变得困难。
事实上，通用I/O是不可撤销的：
一旦你按下了发射核弹头的那个众所周知的按钮，
就再也无法回头了。

以下是在事务中处理I/O的一些选项：

\begin{enumerate}
\item	将事务内的I/O限制为使用内存缓冲区的缓冲I/O。
	这些缓冲区可以像包含任何其他内存位置一样被包含在事务中。
	这似乎是首选机制，而且在许多常见的流I/O和
	大容量存储I/O场景中确实工作得很好。
	然而，在多个面向记录的输出流从多个进程合并到单个文件的情况下，
	需要特殊处理，例如使用\co{fopen()}的``a+''选项
	或\co{open()}的\co{O_APPEND}标志时。
	此外，如下一节所述，常见的网络操作无法通过缓冲来处理。
\item	禁止事务内的I/O，使得任何执行I/O操作的尝试都会中止
	包含它的事务（也许还有多个嵌套事务）。
	这种方法似乎是非缓冲I/O的传统TM处理方式，
	但要求TM与容忍I/O的其他同步原语互操作。
\item	禁止事务内的I/O，但借助编译器来强制执行这一禁令。
\item	在任何给定时间只允许一个特殊的
	\emph{不可撤销}事务~\cite{SpearMichaelScott2008InevitableSTM}
	进行，
	从而允许不可撤销事务包含I/O操作。\footnote{
		在早期文献中，不可撤销事务被称为
		\emph{inevitable} 事务。}
	这在一般情况下可行，但严重限制了I/O操作的可扩展性和性能。
	鉴于可扩展性和性能是并行性的首要目标，
	这种方法的通用性似乎有些自我限制。
	更糟糕的是，使用不可撤销性来容忍I/O操作似乎极大地限制了
	手动事务中止操作的使用。\footnote{
		这一困难由Michael Factor指出。
		要理解这个问题，请思考一下当尝试在事务执行了
		不可撤销操作之后中止该事务时，TM应该做什么。}
	最后，如果有一个不可撤销事务正在操作某个给定的数据项，
	那么操作同一数据项的任何其他事务都不能具有非阻塞语义。
\item	创建新的硬件和协议，使I/O操作能够被纳入事务基础设施中。
	在输入操作的情况下，硬件需要正确预测操作的结果，
	并在预测失败时中止事务。
\end{enumerate}

I/O操作是TM众所周知的弱点，而且在事务中支持I/O的问题
是否有合理的通用解决方案尚不清楚，至少如果``合理''
要包含可用的性能和可扩展性的话。
尽管如此，持续对这一问题的关注和投入很可能会带来更多进展。

\subsubsection{RPC操作}
\label{sec:future:RPC Operations}

可以在基于锁的临界区内执行RPC，在持有hazard pointer时，
在sequence-locking读侧临界区内，以及在userspace-RCU读侧
临界区内执行RPC，如果需要的话，甚至可以同时在所有这些情况下执行。
当你尝试从事务内部执行RPC时会发生什么？

如果RPC请求及其响应都要包含在事务中，并且事务的某个部分
依赖于响应返回的结果，那么就不可能使用在缓冲I/O情况下
可以使用的内存缓冲区技巧。
任何尝试采用这种缓冲方法都会使事务死锁，因为请求要到事务
保证成功后才能发送，但事务是否成功可能要到收到响应后才能
知道，如以下示例所示：

\begin{VerbatimN}[samepage=true]
begin_trans();
rpc_request();
i = rpc_response();
a[i]++;
end_trans();
\end{VerbatimN}

事务的内存足迹要在收到RPC响应之后才能确定，而在事务的内存足迹
确定之前，不可能判断事务是否能够被允许提交。
因此，与事务语义一致的唯一操作就是无条件中止事务，
这至少可以说是毫无帮助的。

以下是TM可用的一些选项：

\begin{enumerate}
\item	禁止事务内的RPC，使得任何执行RPC操作的尝试都会中止
	包含它的事务（也许还有多个嵌套事务）。
	或者，借助编译器来强制执行无RPC的事务。
	这种方法确实可行，但需要TM与其他同步原语交互。
\item	在任何给定时间只允许一个特殊的
	不可撤销事务~\cite{SpearMichaelScott2008InevitableSTM}
	进行，从而允许不可撤销事务包含RPC操作。
	这在一般情况下可行，但严重限制了RPC操作的可扩展性和性能。
	鉴于可扩展性和性能是并行性的首要目标，
	这种方法的通用性似乎有些自我限制。
	此外，使用不可撤销事务来允许RPC操作会在RPC操作开始后
	限制手动事务中止操作。
	最后，如果有一个不可撤销事务正在操作某个给定的数据项，
	那么操作同一数据项的任何其他事务都必须具有阻塞语义。
\item	识别在收到RPC响应之前就能确定事务是否成功的特殊情况，
	并在发送RPC请求之前自动将这些事务转换为不可撤销事务。
	当然，如果多个并发事务以这种方式尝试RPC调用，
	可能需要回滚除一个以外的所有事务，
	从而导致性能和可扩展性下降。
	尽管如此，对于以RPC结尾的长时间运行事务，
	这种方法可能仍然有价值。
	这种方法仍然必须限制手动事务中止操作。
\item	识别可以将RPC响应移出事务的特殊情况，
	然后使用类似于缓冲I/O的技术来处理。
\item	将事务基础设施扩展到包含RPC服务器及其客户端。
	理论上这是可能的，正如分布式数据库所展示的那样。
	然而，鉴于基于内存的TM没有慢速磁盘驱动器
	可以用来隐藏这些延迟，分布式数据库技术能否满足
	所需的性能和可扩展性要求尚不清楚。
	当然，鉴于固态硬盘的出现，数据库也很有可能需要
	重新设计其延迟隐藏方法。
\end{enumerate}

如前一节所述，I/O是TM已知的弱点，而RPC只是I/O中
一个特别棘手的情况。

\subsubsection{时间延迟}
\label{sec:future:Time Delays}

与事务外访问交互的一个重要特殊情况涉及事务内的显式时间延迟。
当然，事务内的时间延迟与TM的原子性属性相矛盾，
但这种情况可以说正是弱原子性要处理的问题。
此外，与内存映射I/O的正确交互有时需要精心控制的时序，
而应用程序经常出于各种目的使用时间延迟。
最后，可以在基于锁的临界区内执行时间延迟，
在持有hazard pointer时，在sequence-locking读侧临界区内，
以及在userspace-RCU读侧临界区内执行时间延迟，
如果需要的话，甚至可以同时在所有这些情况下执行。
从竞争或可扩展性的角度来看，这样做可能不明智，
但话说回来，这样做不会引发任何根本性的概念问题。

那么，TM如何处理事务内的时间延迟？

\begin{enumerate}
\item	忽略事务内的时间延迟。
	这看起来很优雅，但像太多其他``优雅''的解决方案一样，
	无法经受遗留代码的第一次考验。
	这样的代码可能在临界区中有重要的时间延迟，
	在被事务化时会失败。
\item	在遇到时间延迟操作时中止事务。
	这很有吸引力，但不幸的是，并不总是能够自动检测到
	时间延迟操作。
	那个紧凑的循环是在执行关键计算，还是仅仅在等待时间流逝？
\item	借助编译器来禁止事务内的时间延迟。
\item	让时间延迟正常执行。
	不幸的是，一些TM实现只在提交时才发布修改，
	这可能会违背时间延迟的目的。
\end{enumerate}

目前还不清楚是否有一个单一的正确答案。
具有弱原子性并在事务内立即发布更改（在中止时回滚这些更改）
的TM实现可能可以很好地使用最后一种方案。
即使在这种情况下，事务另一端的代码（甚至可能是硬件）
也可能需要大幅重新设计以容忍被中止的事务。
这种重新设计的需求将使得将transactional memory
应用于遗留代码变得更加困难。

\subsubsection{持久性}
\label{sec:future:Persistence}

锁原语有许多不同的类型。
一个有趣的区别是持久性，换句话说，锁是否能够独立于使用该锁的
进程的地址空间而存在。

非持久锁包括\co{pthread_mutex_lock()}、
\co{pthread_rwlock_rdlock()}以及大多数内核级锁原语。
如果实例化非持久锁数据结构的内存位置消失，锁也随之消失。
对于\co{pthread_mutex_lock()}的典型使用，这意味着当进程退出时，
其所有的锁都会消失。
这一特性可以被利用来简化程序关闭时的锁清理工作，
但使得不相关的应用程序之间共享锁变得更加困难，
因为这种共享需要应用程序共享内存。

\QuickQuiz{
	但假设一个应用程序在持有恰好位于文件映射内存区域中的
	\co{pthread_mutex_lock()}时退出了呢？
}\QuickQuizAnswer{
	确实，在这种情况下锁会持久存在，这会让其他试图获取
	这个由已不存在的进程持有的锁的进程感到沮丧。
	这就是为什么在使用位于文件映射内存区域中的\co{pthread_mutex}
	对象时需要格外小心。
}\QuickQuizEnd

持久锁有助于避免不相关应用程序之间共享内存的需求。
持久锁API包括flock系列、\co{lockf()}、System V信号量，
或\co{open()}的\co{O_CREAT}标志。
这些持久API可用于保护跨越多个应用程序运行的大规模操作，
而且在\co{O_CREAT}的情况下甚至能在操作系统重启后存活。
如果需要的话，锁甚至可以通过分布式锁管理器和分布式文件系统
跨越多个计算机系统——并在这些计算机系统中任何一个或全部重启后持续存在。

持久锁可以被任何应用程序使用，包括使用多种语言和软件环境编写的应用程序。
事实上，一个持久锁完全可能由一个用C编写的应用程序获取，
并由一个用Python编写的应用程序释放。

如何为TM提供类似的持久功能？

\begin{enumerate}
\item	将持久事务限制在专门设计来支持它们的特殊用途环境中，
	例如SQL。
	鉴于数据库系统几十年的历史，这显然是可行的，
	但不能提供持久锁所提供的同等程度的灵活性。
\item	使用某些存储设备和/或文件系统提供的快照功能。
	不幸的是，这无法处理网络通信，
	也无法处理对不提供快照功能的设备的I/O，例如U盘。
\item	制造一台时光机。
\item	完全避免这个问题，使用现有的持久设施，
	大概要避免在事务内使用这些设施。
\end{enumerate}

当然，它被称为transactional \emph{memory}这一事实
应该让我们三思，因为这个名称本身就与持久事务的概念相矛盾。
不过，将这种可能性作为探测transactional memory固有局限性的
重要测试用例来考虑仍然是值得的。

\subsection{进程修改}
\label{sec:future:Process Modification}

进程不是永恒的：
它们被创建和销毁，它们的内存映射被修改，
它们与动态库链接，它们被调试。
这些小节探讨transactional memory如何处理不断变化的执行环境。

\subsubsection{多线程事务}
\label{sec:future:Multithreaded Transactions}

在持有锁的同时创建进程和线程是完全合法的，
或者更确切地说，在持有hazard pointer时，在sequence-locking
读侧临界区内，以及在userspace-RCU读侧临界区内也是如此，
如果需要的话，甚至可以同时在所有这些情况下执行。
这不仅是合法的，而且非常简单，从以下代码片段可以看出：

\begin{VerbatimN}
pthread_mutex_lock(...);
for (i = 0; i < ncpus; i++)
	pthread_create(&tid[i], ...);
for (i = 0; i < ncpus; i++)
	pthread_join(tid[i], ...);
pthread_mutex_unlock(...);
\end{VerbatimN}

这段伪代码使用\co{pthread_create()}为每个CPU生成一个线程，
然后使用\co{pthread_join()}等待每个线程完成，
全部在\co{pthread_mutex_lock()}的保护下进行。
其效果是并行执行一个基于锁的临界区，
使用\co{fork()}和\co{wait()}也可以获得类似的效果。
当然，临界区需要足够大才能证明线程创建开销的合理性，
但在生产软件中有许多大型临界区的例子。

TM如何处理事务内的线程创建？

\begin{enumerate}
\item	声明\co{pthread_create()}在事务内是非法的，
	最好是通过中止事务来实现。
	或者，借助编译器来强制执行无\co{pthread_create()}的事务。
\item	允许在事务内执行\co{pthread_create()}，
	但只有父线程被视为事务的一部分。
	这种方法似乎与现有的和设想的TM实现合理兼容，
	但对于不知情者来说似乎是一个陷阱。
	这种方法引发了更多问题，例如如何处理冲突的子线程访问。
\item	将\co{pthread_create()}转换为函数调用。
	这种方法也是一个有吸引力的陷阱，因为它无法处理
	子线程之间相互通信这一并不罕见的情况。
	此外，它不允许事务主体的并发执行。
\item	将事务扩展为覆盖父线程和所有子线程。
	这种方法引发了关于冲突访问性质的有趣问题，
	因为父线程和子线程大概被允许相互冲突，
	但不能与其他线程冲突。
	它还引发了如果父线程在提交事务之前不等待其子线程
	会发生什么的有趣问题。
	更有趣的是，如果父线程根据参与事务的变量值
	有条件地执行\co{pthread_join()}会怎样？
	在锁的情况下，这些问题的答案相当直接。
	TM的答案留给读者作为练习。
\end{enumerate}

鉴于事务的并行执行在数据库世界中是司空见惯的，
当前的TM提案不支持这一点也许令人惊讶。
另一方面，上面的例子是锁的一种相当复杂的用法，
在简单的教科书示例中通常看不到，
所以它的缺失也许是可以预期的。
话虽如此，一些研究人员正在使用事务来自动并行化
代码~\cite{ArunRaman2010MultithreadedTransactions}，
据传其他TM研究人员正在研究事务内的fork/join并行性，
因此这个话题也许很快就会得到更全面的讨论。

\subsubsection{\tco{exec()}系统调用}
\label{sec:future:The exec System Call}

可以在基于锁的临界区内执行\co{exec()}系统调用，
在持有hazard pointer时，在sequence-locking读侧临界区内，
以及在userspace-RCU读侧临界区内执行，
如果需要的话，甚至可以同时在所有这些情况下执行。
确切的语义取决于原语的类型。

对于非持久原语（包括\co{pthread_mutex_lock()}、
\co{pthread_rwlock_rdlock()}和userspace RCU），
如果\co{exec()}成功，整个地址空间消失，
连同任何正在持有的锁。
当然，如果\co{exec()}失败，地址空间仍然存在，
因此任何相关的锁也仍然存在。
也许有点奇怪，但定义明确。

另一方面，持久原语（包括flock系列、\co{lockf()}、
System V信号量和\co{open()}的\co{O_CREAT}标志）
无论\co{exec()}成功还是失败都会继续存在，
因此被\co{exec()}执行的程序完全可以释放它们。

\QuickQuiz{
	那么由\co{mmap()}内存区域中的数据结构表示的
	非持久原语呢？
	当在这种原语的临界区内执行\co{exec()}时会发生什么？
}\QuickQuizAnswer{
	如果被\co{exec()}执行的程序映射了相同的内存区域，
	那么该程序原则上可以简单地释放锁。
	从软件工程的角度来看这种方法是否合理，
	留给读者作为练习。
}\QuickQuizEnd

当你尝试从事务内部执行\co{exec()}系统调用时会发生什么？

\begin{enumerate}
\item	禁止事务内的\co{exec()}，使得包含它的事务在遇到\co{exec()}时中止。
	这定义明确，但显然需要非TM同步原语来配合\co{exec()}使用。
\item	禁止事务内的\co{exec()}，由编译器强制执行此禁令。
	有一个C++ TM草案规范采用了这种方法，
	允许函数使用\co{transaction_safe}和\co{transaction_unsafe}
	属性进行修饰。\footnote{
		感谢Mark Moir向我指出了这个规范，
		也感谢Michael Wong此前向我指出了更早的修订版本。}
	这种方法相比在运行时中止事务有一些优势，
	但同样需要非TM同步原语来配合\co{exec()}使用。
	一个缺点是需要为大量库函数添加\co{transaction_safe}和
	\co{transaction_unsafe}属性。
\item	以类似于非持久锁原语的方式对待事务，
	使得如果\co{exec()}失败则事务继续存在，
	如果\co{exec()}成功则静默提交。
	只有部分受事务影响的变量驻留在\co{mmap()}内存中
	（因此可以在\co{exec()}系统调用成功后存活）的情况
	留给读者作为练习。
\item	如果\co{exec()}系统调用本来会成功，
	则中止事务（和\co{exec()}系统调用），
	但如果\co{exec()}系统调用本来会失败，
	则允许事务继续。
	这在某种意义上是``正确''的方法，
	但需要大量工作却得到相当不令人满意的结果。
\end{enumerate}

\co{exec()}系统调用也许是TM普遍适用性障碍中最奇特的例子，
因为目前还不完全清楚哪种方法有意义，
有些人可能会认为这仅仅是与\co{exec()}的现实交互
中的困境的反映。
话虽如此，禁止事务内\co{exec()}的两个选项
也许是这组选项中最合乎逻辑的。

类似的问题也围绕着\co{exit()}和\co{kill()}系统调用，
以及会退出事务的\co{longjmp()}或异常。
（\co{longjmp()}或异常是从哪里来的？）

\subsubsection{动态链接和加载}
\label{sec:future:Dynamic Linking and Loading}

基于锁的临界区、持有hazard pointer的代码、
sequence-locking读侧临界区和userspace-RCU读侧临界区
可以（单独或组合地）合法地包含调用动态链接和加载的函数的代码，
包括C/C++共享库和Java类库。
当然，这些库中包含的代码在编译时根据定义是不可知的。
那么，如果在事务内调用一个动态加载的函数会发生什么？

这个问题有两个部分：
\begin{enumerate*}[(a)]
\item 如何在事务内动态链接和加载一个函数，以及
\item 如何处理该函数内代码的不可知性质？
\end{enumerate*}
公平地说，问题(b)至少在理论上也给锁和userspace-RCU带来了一些挑战。
例如，动态链接的函数可能会为锁引入死锁（\IX{deadlock}），
或者可能（错误地）在userspace-RCU读侧临界区中引入
\IX{quiescent state}（静止状态）。
区别在于，虽然锁和userspace-RCU临界区中允许的操作类别已被充分理解，
但在TM的情况下似乎仍然存在相当大的不确定性。
事实上，不同的TM实现似乎有不同的限制。

那么TM如何处理动态链接和加载的库函数呢？
对于部分(a)即代码的实际加载，选项包括以下内容：

\begin{enumerate}
\item	以类似于缺页错误的方式处理动态链接和加载，
	使得函数被加载和链接，可能在此过程中中止事务。
	如果事务被中止，重试将发现函数已经存在，
	因此可以预期事务正常进行。
\item	禁止在事务内动态链接和加载函数。
\end{enumerate}

对于部分(b)即无法检测尚未加载的函数中TM不友好的操作，
可能的选项包括以下内容：

\begin{enumerate}
\item	直接执行代码：
	如果函数中有任何TM不友好的操作，则简单地中止事务。
	不幸的是，这种方法使编译器无法确定一组给定的事务
	是否可以安全地组合。
	无论如何都允许可组合性的一种方式是不可撤销事务，
	然而当前的实现在任何给定时间只允许一个不可撤销事务进行，
	这可能严重限制性能和可扩展性。
	不可撤销事务也会限制手动事务中止操作的使用。
	最后，如果有一个不可撤销事务正在操作某个给定的数据项，
	操作同一数据项的任何其他事务都不能具有非阻塞语义。
\item	修饰函数声明以指示哪些函数是TM友好的。
	这些修饰可以由编译器的类型系统来强制执行。
	当然，对于许多语言来说，这需要提出、标准化和实现语言扩展，
	带有相应的时间延迟，以及对大量原本无关的库函数
	进行相应的修饰。
	话虽如此，标准化工作已经在
	进行中~\cite{Ali-Reza-Adl-Tabatabai2009CppTM,Ali-Reza-Adl-Tabatabai2012CppTM}。
\item	如上所述，禁止在事务内动态链接和加载函数。
\end{enumerate}

I/O操作当然是TM已知的弱点，而动态链接和加载
可以被视为I/O的又一个特殊情况。
尽管如此，TM的支持者必须要么解决这个问题，
要么接受TM只是并行程序员工具箱中几个工具之一这样的现实。
（公平地说，许多TM支持者早已接受了一个不仅仅包含TM的世界。）

\subsubsection{内存映射操作}
\label{sec:future:Memory-Mapping Operations}

在基于锁的临界区内执行内存映射操作（包括\co{mmap()}、
\co{shmat()}和\co{munmap()}~\cite{TheOpenGroup1997SUS}）
是完全合法的，在持有hazard pointer时、
在sequence-locking读侧临界区内、以及在userspace-RCU读侧
临界区内也是如此，如果需要的话甚至可以同时在所有这些情况下执行。
当你尝试从事务内部执行这样的操作时会发生什么？
更关键的是，如果被重映射的内存区域包含参与当前线程事务的
某些变量会怎样？
如果这个内存区域包含参与其他线程事务的变量又会怎样？

不需要考虑TM系统的元数据被重映射的情况，
因为大多数锁原语都不定义重映射其锁变量的结果。

以下是一些TM内存映射选项：

\begin{enumerate}
\item	事务内的内存重映射是非法的，将导致所有包含它的事务被中止。
	这确实在某种程度上简化了问题，但也要求TM与那些在其临界区内
	容忍重映射的同步原语互操作。
\item	事务内的内存重映射是非法的，由编译器强制执行此禁令。
\item	事务内的内存映射是合法的，但会中止所有在被映射区域中
	拥有变量的其他事务。
\item	事务内的内存映射是合法的，但如果被映射的区域与当前事务的
	内存足迹重叠，映射操作将失败。
\item	所有内存映射操作，无论是在事务内还是事务外，
	都检查被映射区域是否与系统中所有事务的内存足迹冲突。
	如果有重叠，则内存映射操作失败。
\item	与系统中任何事务的内存足迹重叠的内存映射操作的效果
	由TM冲突管理器决定，它可以动态确定是让内存映射操作失败
	还是中止任何冲突的事务。
\end{enumerate}

值得注意的是，\co{munmap()}会使相关内存区域处于未映射状态，
这可能有额外的有趣含义。\footnote{
	映射和取消映射之间的这种差异由Josh Triplett指出。}

\subsubsection{调试}
\label{sec:future:Debugging}

通常的调试操作如断点在基于锁的临界区内和userspace-RCU读侧
临界区内正常工作。
然而，在最初的transactional-memory硬件
实现~\cite{DaveDice2009ASPLOSRockHTM}中，事务内的异常
将中止该事务，这反过来意味着断点会中止所有包含它的事务。

那么如何调试事务？

\begin{enumerate}
\item	在包含断点的事务中使用软件仿真技术。
	当然，每当在任何事务的范围内设置断点时，
	可能需要仿真所有事务。
	如果运行时系统无法确定给定的断点是否在事务的范围内，
	那么为了安全起见，可能需要仿真所有事务。
	然而，这种方法可能会带来显著的开销，
	这反过来可能会掩盖正在追踪的bug。
\item	仅使用能够处理断点异常的硬件TM实现。
	不幸的是，截至本文撰写时（2021年3月），所有此类
	实现都是研究原型。
\item	仅使用软件TM实现，它们（非常粗略地说）
	比较简单的硬件TM实现对异常更具容忍性。
	当然，软件TM的开销往往比硬件TM更高，
	因此这种方法可能并非在所有情况下都可接受。
\item	更仔细地编程，以便首先避免在事务中出现bug。
	一旦你弄清楚如何做到这一点，请务必让大家都知道这个秘诀！
\end{enumerate}

有一些理由相信transactional memory将带来与其他同步机制相比的
\IX{productivity}（生产力）改进，
但这些改进如果传统调试技术不能应用于事务的话，
似乎很容易被丧失。
如果transactional memory要被新手用于大型事务，
这一点似乎尤其正确。
相比之下，硬派的``顶尖''程序员也许能够不用这些调试辅助工具，
尤其是对于小型事务。

因此，如果transactional memory要兑现其对新手程序员的生产力承诺，
调试问题确实需要得到解决。

\subsection{同步}
\label{sec:future:Synchronization}

如果transactional memory有朝一日证明它能够满足所有人的所有需求，
那么它就不需要与任何其他同步机制交互。
在那之前，它需要与能做它不能做之事的同步机制配合工作，
或者与在给定情况下更自然地工作的同步机制配合。
以下各节概述了该领域当前的挑战。

\subsubsection{锁}
\label{sec:future:Locking}

在持有其他锁的同时获取锁是很常见的做法，这工作得很好，
至少只要使用了通常众所周知的软件工程技术来避免死锁。
从RCU读侧临界区内获取锁也不罕见，
这减轻了死锁方面的担忧，因为RCU读侧原语
不能参与基于锁的死锁循环。
在持有hazard pointer和sequence-lock读侧临界区内也可以获取锁。
但当你尝试从事务内部获取锁时会发生什么？

理论上，答案很简单：
只需将表示锁的数据结构作为事务的一部分进行操作，
一切就会完美运作。
在实践中，根据TM系统的实现细节，可能会出现许多
不明显的复杂情况~\cite{Volos2008TRANSACT}。
这些复杂情况可以解决，但代价是在事务外获取锁的开销增加45\pct，
在事务内获取锁的开销增加300\pct。
虽然这些开销对于包含少量锁的事务性程序可能是可以接受的，
但对于希望偶尔使用事务的生产级基于锁的程序来说，
这些开销往往是完全不可接受的。

以下是TM可用的一些选项：

\begin{enumerate}
\item	仅使用对锁友好的TM实现~\cite{DaveDice2006DISC}。
	不幸的是，对锁不友好的实现具有一些有吸引力的属性，
	包括成功事务的低开销和能够容纳极大事务的能力。
\item	在将TM引入基于锁的程序时仅``小规模''使用TM，
	从而适应对锁友好的TM实现的限制。
\item	完全搁置基于锁的遗留系统，用事务重新实现一切。
	这种方法不缺少倡导者，但这要求
	本系列中描述的所有问题都得到解决。
	在解决这些问题所需的时间内，竞争的同步机制
	当然也有机会得到改进。
\item	严格将TM用作基于锁系统中的优化，正如
	TxLinux~\cite{ChistopherJRossbach2007a}小组和
	大量transactional lock elision
	项目~\cite{MartinPohlack2011HTM2TLE,Kleen:2014:SEL:2566590.2576793,PascalFelber2016rwlockElision,SeongJaePark2020HTMRCUlock}所做的那样。
	这种方法似乎是合理的，但锁的设计约束
	（如需要避免死锁）仍然完全保留。
\item	努力减少施加在锁原语上的开销。
\end{enumerate}

TM和锁之间可能存在接口问题这一事实令许多人感到惊讶，
这突显了在实际生产软件中试用新机制和原语的必要性。
幸运的是，开源的出现意味着大量这样的软件现在对所有人
（包括研究人员）都是免费可用的。

\subsubsection{读写锁}
\label{sec:future:Reader-Writer Locking}

在持有其他锁的同时以读模式获取读写锁（reader-writer lock）是很常见的做法，
这直接就能工作，至少只要使用了通常众所周知的软件工程技术
来避免死锁。
从RCU读侧临界区内以读模式获取读写锁也能工作，
而且这样做减轻了死锁方面的担忧，
因为RCU读侧原语不能参与基于锁的死锁循环。
在持有hazard pointer和sequence-lock读侧临界区内也可以获取锁。
但当你尝试从事务内部以读模式获取读写锁时会发生什么？

不幸的是，在事务中以直接的方式读获取传统的基于计数器的读写锁
会违背读写锁的目的。
要理解这一点，考虑一对事务同时尝试读获取同一个读写锁。
因为读获取涉及修改读写锁的数据结构，
会产生冲突，这将回滚两个事务中的一个。
这种行为与读写锁允许并发读者的目标完全不一致。

以下是TM可用的一些选项：

\begin{enumerate}
\item	使用per-CPU或per-thread读写
	锁~\cite{WilsonCHsieh92a}，它允许给定的CPU
	（或线程）在读获取锁时仅操作本地数据。
	这将避免两个事务同时读获取锁之间的冲突，
	允许两者如预期那样继续进行。
	不幸的是，(1)~per-CPU/thread锁的写获取开销可能极高，
	(2)~per-CPU/thread锁的内存开销可能令人望而却步，
	(3)~这种转换仅在你能访问相关源代码时才可用。
	其他更新的可扩展读写锁~\cite{YossiLev2009SNZIrwlock}
	可能避免部分或全部这些问题。
\item	在将TM引入基于锁的程序时仅``小规模''使用TM，
	从而避免从事务内部读获取读写锁。
\item	完全搁置基于锁的遗留系统，用事务重新实现一切。
	这种方法不缺少倡导者，但这要求
	本系列中描述的\emph{所有}问题都得到解决。
	在解决这些问题所需的时间内，竞争的同步机制
	当然也有机会得到改进。
\item	严格将TM用作基于锁系统中的优化，正如
	TxLinux~\cite{ChistopherJRossbach2007a}小组所做的，
	以及更近期使用TM来省略读写锁的
	工作~\cite{PascalFelber2016rwlockElision}所做的那样。
	这种方法似乎是合理的，至少在\Power{8}
	CPU~\cite{Le:2015:TMS:3266491.3266500}上是这样，
	但锁的设计约束（如需要避免死锁）仍然完全保留。
\end{enumerate}

当然，将TM与读写锁结合可能还有其他不明显的问题，
就像与排他锁结合时实际上确实出现了问题一样。

\subsubsection{延迟回收}
\label{sec:future:Deferred Reclamation}

本节主要关注RCU。
将TM与其他延迟回收机制（如\IXalt{reference counters}{reference count}
和hazard pointers）结合时也会出现类似的问题和可能的解决方案。
在下面的文本中，已知的差异会被特别指出。

引用计数（reference counting）、hazard pointers和RCU都被大量使用，如
\cref{sec:defer:RCU Related Work,sec:defer:Which to Choose? (Production Use)}
中所述。
这意味着任何选择不克服本节中指出的每个挑战的TM实现
都需要与所有这些同步机制干净而高效地互操作。

德克萨斯大学奥斯汀分校的TxLinux小组似乎是
接受RCU/TM互操作性挑战的
小组~\cite{ChistopherJRossbach2007a}。
因为他们将TM应用于使用RCU的Linux 2.6内核，
所以他们别无选择，只能集成TM和RCU，用TM取代锁来进行RCU更新。
不幸的是，虽然论文确实指出RCU实现的锁
（例如\co{rcu_ctrlblk.lock}）被转换为事务，
但对于那些由基于RCU的更新使用的锁
（例如\co{dcache_lock}）做了什么却没有提及。

更近期，\ppl{Dimitrios}{Siakavaras}等人将HTM和RCU应用于搜索
树~\cite{Siakavaras2017CombiningHA,DimitriosSiakavaras2020RCU-HTM-B+Trees}，
\ppl{Christina}{Giannoula}等人使用HTM和RCU来对图进行
着色~\cite{ChristinaGiannoula2018HTM-RCU-graphcoloring}，
而\ppl{SeongJae}{Park}等人使用HTM和RCU来优化\IXacr{numa}系统上的
高竞争锁~\cite{SeongJaePark2020HTMRCUlock}。

重要的是要注意，RCU允许读者和更新者并发运行，
进一步允许RCU读者访问正在被更新的数据。
当然，RCU的这一特性，无论其性能、可扩展性和
实时响应优势如何，都与TM的底层原子性属性相矛盾，
尽管\Power{8} CPU系列的挂起事务
设施~\cite{Le:2015:TMS:3266491.3266500}使其成为这一规则的例外。

那么基于TM的更新应该如何与并发RCU读者交互？
一些可能的方式如下：

\begin{enumerate}
\item	RCU读者中止并发冲突的TM更新。
	这实际上是TxLinux项目采用的方法。
	这种方法确实保留了RCU语义，也保留了RCU的读侧性能、
	可扩展性和实时响应特性，但它有一个不幸的副作用，
	即不必要地中止冲突的更新。
	在最坏的情况下，一长串RCU读者可能饿死所有更新者，
	这在理论上可能导致系统挂起。
	此外，不是所有TM实现都提供实现此方法所需的强原子性，
	而且有充分的理由。
\item	与冲突TM更新并发运行的RCU读者从任何冲突的RCU加载中
	获取旧的（事务前的）值。
	这保留了RCU语义和性能，也防止了RCU更新饥饿（\IX{starvation}）。
	然而，不是所有TM实现都能及时提供对已被进行中事务
	暂时更新的变量旧值的访问。
	特别是，在日志中维护旧值的基于日志的TM实现
	（因此提供了出色的TM提交性能）不太可能喜欢这种方法。
	也许可以利用\co{rcu_dereference()}原语来允许RCU
	在更广泛的TM实现中访问旧值，
	尽管性能可能仍然是个问题。
	尽管如此，有流行的TM实现已经以这种方式与RCU
	集成~\cite{DonaldEPorter2007TRANSACT,PhilHoward2011RCUTMRBTree,
	PhilipWHoward2013RCUrbtree}。
\item	如果RCU读者执行了与进行中事务冲突的访问，
	那么该RCU访问将被延迟，直到冲突事务提交或中止。
	这种方法保留了RCU语义，但不保留RCU的性能或实时响应，
	特别是在存在长时间运行事务的情况下。
	此外，不是所有TM实现都能够延迟冲突的访问。
	尽管如此，对于仅支持小型事务的硬件TM实现，
	这种方法似乎非常合理。
\item	RCU读者被转换为事务。
	这种方法基本上保证了RCU与任何TM实现兼容，
	但它也将TM的回滚强加于RCU读侧临界区，
	破坏了RCU的实时响应保证，也降低了RCU的读侧性能。
	此外，在任何RCU读侧临界区包含相关TM实现无法处理的
	操作的情况下，这种方法是不可行的。
	这种方法更难应用于hazard pointers和引用计数，
	因为它们没有将读者作为代码段的明确定义。
\item	RCU的许多更新侧用法修改单个指针来发布新的数据结构。
	在其中一些情况下，只要事务尊重内存排序并且回滚过程
	使用\co{call_rcu()}来释放相应的结构，
	就可以安全地允许RCU看到随后被回滚的事务性指针更新。
	不幸的是，不是所有TM实现都尊重事务内的内存屏障。
	显然，人们的想法是，因为事务应该是原子的，
	事务内访问的顺序不应该有影响。
\item	禁止在RCU更新中使用TM。
	这保证可以工作，但限制了TM的使用。
\end{enumerate}

可能还会发现更多的方法，特别是考虑到用户级RCU和
hazard-pointer实现的出现。\footnote{
	感谢TxLinux小组、Maged Michael和Josh Triplett
	提出了上述多种替代方案。}
有趣的是，许多性能和可扩展性较好的STM实现在内部使用了
类RCU技术~\cite{UCAM-CL-TR-579,KeirFraser2007withoutLocks,Gu:2019:PSE:3358807.3358885,Kim:2019:MSR:3297858.3304040}。

\QuickQuiz{
	MV-RLU看起来相当不错！
	它是不是完胜RCU？
}\QuickQuizAnswer{
	快速阅读摘要可能会给人这种印象，
	但更仔细的读者会注意到最后一句中的``对于广泛的工作负载''
	这一短语。
	事实证明这个短语相当重要：

	\begin{enumerate}
	\item	他们的RCU评估使用了同步宽限期，这不必要地
		限制了更新速度，正如他们的第6.2.1节所述。
		参见本书的\cref{fig:datastruct:Read-Side RCU-Protected Hash-Table Performance For Schroedinger's Zoo in the Presence of Updates}
		\cpageref{fig:datastruct:Read-Side RCU-Protected Hash-Table Performance For Schroedinger's Zoo in the Presence of Updates}，
		可以看到历史悠久的异步\co{call_rcu()}原语
		使RCU能够在大量更新者的情况下表现和扩展得相当好。
		此外，在他们论文的第3.7节中，作者承认
		异步宽限期对MV-RLU的可扩展性很重要。
		公平的比较也应该允许RCU享受异步的好处。
	\item	他们使用了一个调优不佳的含10,000个元素的
		1,000桶哈希表。
		此外，他们的448个硬件线程需要远多于1,000个桶
		才能避免他们正确指出的限制RCU在其基准测试中
		性能的锁竞争（lock contention）。
		有用的比较应该使用适当调优的哈希表。
	\item	他们的RCU哈希表使用了per-bucket锁，
		他们将其称为瓶颈，鉴于长哈希链和
		桶与线程的小比率，这并不令人惊讶。
		他们的一些竞争机制改为使用无锁（lock-free）技术，
		从而避免了per-bucket锁的瓶颈，
		愤世嫉俗者可能会声称这揭示了作者在其他方面
		令人费解地选择调优不佳哈希表的原因。
		作者图4中间行的第一张图显示了RCU在不受
		人为瓶颈束缚时能达到的效果，
		同一行第二张图的前半部分也是如此。
	\item	他们的链表操作允许RLU对列表中不同元素进行并发修改，
		而RCU被迫串行化更新。
		同样，RCU一直与无锁更新者配合得很好，
		这一事实已在作者引用的学术文献中
		阐述~\cite{MathieuDesnoyers2012URCU}。
		公平的比较应该对RCU使用与MV-RLU相同风格的更新。
	\item	作者没有考虑将RCU和sequence locking结合使用，
		Linux内核中使用这种组合来给读者
		提供多指针更新的一致视图。
	\item	作者没有考虑基于RCU的Issaquah
		Challenge~\cite{PaulEMcKenney2016IssaquahCPPCON}解决方案，
		该方案也给读者提供多指针更新的一致视图，
		尽管对``一致''的定义较弱。
	\end{enumerate}

	令人惊讶的是，这篇论文的匿名审稿人没有要求
	对MV-RLU和RCU进行同类比较。
	尽管如此，应该祝贺作者撰写了一篇学术论文，
	其中展示了良好可扩展性与强读侧一致性相结合的罕见例子。
	也应该祝贺他们克服了传统学术界对异步宽限期的偏见，
	这极大地帮助了他们的可扩展性。

	有趣的是，RLU和RCU采用不同的方法来避免
	\ppl{Hagit}{Attiya}等人指出的STM固有
	限制~\cite{Attiya:2009:STMReadOnlyLimits}。
	RCU避免提供严格的可串行化，RLU避免提供
	不可见的只读事务，两者因此都避免了这些限制。
}\QuickQuizEnd

\subsubsection{事务外访问}
\label{sec:future:Extra-Transactional Accesses}

在基于锁的临界区内，操作在该锁的临界区外被并发访问
甚至修改的变量是完全合法的，一个常见的例子是统计计数器（statistical counter）。
在RCU读侧临界区内也可以做同样的事情，
而且这实际上是常见情况。

鉴于生产数据库系统中普遍存在的所谓``脏读''等机制，
TM的支持者认真关注事务外访问就不足为奇了，
弱原子性~\cite{Blundell2006TMdeadlock}的概念就是一个例子。

以下是一些事务外访问选项：

\begin{enumerate}
\item	由事务外访问引起的冲突总是中止事务。
	这就是强原子性。
\item	由事务外访问引起的冲突被忽略，
	因此只有事务之间的冲突才能中止事务。
	这就是弱原子性。
\item	允许事务在特殊情况下执行非事务操作，
	例如在分配内存或与基于锁的临界区交互时。
\item	产生允许某些操作（例如加法）由多个事务
	在单个变量上并发执行的硬件扩展。
\item	向transactional memory引入弱语义。
	一种方法是\cref{sec:future:Deferred Reclamation}中描述的
	与RCU的结合，而Gramoli和Guerraoui
	调查了许多其他弱事务
	方法~\cite{Gramoli:2014:DTP:2541883.2541900}，例如，
	将大型``弹性''事务限制性地分割为较小的事务，
	从而降低冲突概率（尽管性能和可扩展性不温不火）。
	也许进一步的经验将表明事务外访问的某些用途
	可以被弱事务所取代。
\end{enumerate}

看起来事务是在真空中构想出来的，
不需要与任何其他同步机制交互。
如果是这样，当事务与非事务访问结合时出现大量混乱和复杂性
就不足为奇了。
但除非事务被限制在对孤立数据结构的小型更新上，
或者被限制在不与大量现有并行代码交互的新程序上，
否则如果事务要在近期产生大规模实际影响，
就绝对必须进行这种结合。

\subsection{其他事务}
\label{sec:future:Other Transactions}
% If we ever get epigraphs down at this level, perhaps the old
% "Hell is other people" quote by Jean-Paul Sartre.  His clarification
% also applies, which can be stated as "Heaven is also other people".

因为事务对彼此应该表现为原子的，
当多个事务同时尝试访问同一个变量时必须采取一些措施。
如果所有事务都在读取该变量而没有一个在更新它，
那么可以``什么都不做''，否则事务必须至少部分串行化。
这种串行化通常通过中止和回滚那些事务的某个子集来实现。

如何选择中止和回滚哪些事务？

关于这个问题已经有大量的工作。
答案是``使用竞争管理器''，但这只是将问题转移到了
``竞争管理器应该做什么？''
对竞争管理器的全面论述超出了本节的范围。
因此本节将限于该领域的几篇开创性论文。

Herlihy等人~\cite{HerlihyLMS03}
让写事务无条件地中止读事务，
这有简单性的优点，但可能导致读者饥饿。
这篇论文还引入了早期释放的概念，
它可以允许一次长遍历与小更新并发进行。
从这个角度来看，可以将RCU读者描述为``即时释放''，
但这是一种使RCU读者以完全不同于STM方式行为的特性。

Hammond等人~\cite{LanceHammond2004a}
建议使用小事务，这确实降低了竞争管理器的复杂性，
但也未能兑现STM自动并发的全部承诺。
这篇论文还引入了注解来标识不同类型的事务。

Scherer和Scott~\cite{WilliamNSchererIII2005}
描述了许多STM竞争管理策略。
这篇论文为事务外访问定义了``可见''和``不可见''的读者，
并讨论了许多竞争管理策略，包括：

\begin{enumerate}
\item	指数退避，他们得出结论认为需要手动调优。
\item	优先中止完成工作量最少的事务，
	其中之前被中止的尝试完成事务时所做的工作
	会计入该事务下次重试的工作量。
	这有很好的公平性属性，但可能导致大型事务
	长时间占用系统。
\item	优先中止完成工作量最少的事务，
	但也将等待给定事务完成的任何其他事务所做的工作
	计入该事务。
	这在组合中加入了一种优先级继承的形式。
\item	跟踪给定事务被另一个事务中止的次数，
	作为促进某种公平性概念的另一种方式。
\item	使用事务开始时间戳，以便优先中止最近开始的事务。
\item	使用事务开始时间戳，但也添加给定事务上次运行时间的指示，
	以防止被抢占的事务中止在更高优先级进程中运行的其他事务。
\end{enumerate}

这篇2005年的论文因此早期指出了竞争管理中固有的复杂性。

Porter和Witchel~\cite{DonaldEPorter2007TRANSACT,Ramadan:2008:DTM:1521747.1521799}
考虑放宽STM一致性要求以简化竞争管理。
他们指出RCU是更新会导致并发读取访问数据旧版本的一种情况。

竞争管理在HTM中也是一个问题，将在
\cref{sec:future:Conflict Handling,sec:future:Aborts and Rollbacks}
中讨论。

\QuickQuiz{
	为什么不与时俱进，将机器学习应用于竞争管理？
}\QuickQuizAnswer{
	许多事务的开销在亚微秒级别，因此在机器学习上消耗
	比通过更高效地执行事务所节省的更多CPU是太容易的事情。
	尽管如此，机器学习可能在将竞争管理器调优到
	特定工作负载方面发挥作用。
	当然，不需要调优会更好，但有时``更好''是不可实现的。
}\QuickQuizEnd

\subsection{案例研究：
	    Sequence Locking}
\label{sec:future:Case Study: Sequence Locking}

\cref{sec:defer:Sequence Locks}中描述的Sequence locking
有时被认为是STM的一种简单形式。
Sequence-locking读侧临界区通常限于执行加载操作，
每当与sequence-locking更新者发生冲突时就会被重试。
至关重要的是，sequence-locking读侧临界区之间不会相互冲突。
因此，考虑sequence locking如何处理前面各节中指出的
STM挑战是很有启发性的。

Sequence locking通过简单地在每次通过其读侧临界区时
执行TM不友好的操作来处理与外部世界的交互
（\cref{sec:future:Outside World}）。
例如，如果执行I/O写入的临界区被重试了三次，
它将执行四次该I/O写入，一次是原始通过临界区的执行，
三次是每次重试各一次。

在sequence-locking读侧临界区内执行I/O操作
可能看起来有些非常规。
然而，出于历史记录、调试或审计目的
向日志进行I/O写入（\cref{sec:future:I/O Operations}）
是非常有意义的。
RPC操作（\cref{sec:future:RPC Operations}）可以收集所需数据。
时间延迟（\cref{sec:future:Time Delays}）可以为调试目的
提供有价值的服务，例如测试重试代码路径。
即使是持久操作（\cref{sec:future:Persistence}）也可能有用，
例如执行跨应用程序的互斥。
或者，就像STM一样，sequence locking的用户可以选择
缓冲I/O并在读侧临界区结束后立即执行它。

Sequence locking以相同的方式处理进程修改
（\cref{sec:future:Process Modification}），
即在其读侧临界区的每次执行中简单地执行进程修改操作。

在临界区内创建进程和线程
（\cref{sec:future:Multithreaded Transactions}）
正常运作，尽管应该注意防止经常重试的临界区成为``fork炸弹''。
\co{exec()}系统调用（\cref{sec:future:The exec System Call}）
有完全合理的（虽然相当激进的）语义。
对动态库中函数的调用
（\cref{sec:future:Dynamic Linking and Loading}），
例如Linux内核可加载模块，与对内置函数的调用一样正常工作。
当然，如果这样的调用导致模块加载，
这可能导致临界区第一次执行时的高重试概率。
重映射内存（\cref{sec:future:Memory-Mapping Operations}）
行为直观，至少假设组成临界区的代码和sequence lock本身
都没有被重映射。
请注意，这些相同的限制也适用于普通锁。
调试器（\cref{sec:future:Debugging}）
在sequence-locking读侧临界区内正常工作，
就像在普通锁的临界区内一样。

同步原语（\cref{sec:future:Synchronization}）
可以在sequence-locking读侧临界区内自由使用。
锁（\cref{sec:future:Locking}）、
读写锁（\cref{sec:future:Reader-Writer Locking}）、
延迟回收（\cref{sec:future:Deferred Reclamation}）
以及事务外访问
（\cref{sec:future:Extra-Transactional Accesses}）
都可以在sequence-locking读侧临界区内使用，
并且都给出预期的结果。

Sequence locking使用一个简单的竞争管理器
（\cref{sec:future:Other Transactions}）。
Sequence-locking读者在有并发更新者的情况下总是被回滚
（正如Herlihy等人~\cite{HerlihyLMS03}所做的），
sequence-locking更新者通常将其竞争管理委托给某种形式的互斥，
而sequence-locking读者从不冲突，
因此从不需要管理其不存在的竞争。
这种粗略直接的方法意味着sequence-locking更新者可能饿死读者，
因此sequence locking的用户必须注意避免频繁更新。

简而言之，sequence locking相当好地处理了这些STM挑战。
这可能是sequence locking在生产中被大量使用的原因之一。
然而，sequence locking的一个关键使能因素是
将许多这些挑战推回给用户。
这种方法是合理的，因为sequence locking没有成为
统治所有同步机制的唯一机制的野心。
因此，也许广泛采用STM的最大障碍是
其支持者对完全通用性的过度渴望。

\subsection{讨论}
\label{sec:future:Discussion}

TM普遍采用的障碍导致以下结论：

\begin{enumerate}
\item	TM的一个有趣属性是事务可能被回滚和重试。
	这一属性是TM在不可逆操作方面困难的根源，
	包括非缓冲I/O、RPC、内存映射操作、
	时间延迟和\co{exec()}系统调用。
	这一属性还有一个不幸的后果，即引入了失败可能性
	中固有的所有复杂性，而且通常以开发者可见的方式体现。
\item	TM的另一个有趣属性，由Shpeisman
	等人~\cite{TatianaShpeisman2009CppTM}指出，是TM将同步
	与它保护的数据交织在一起。
	这一属性是TM在I/O、内存映射操作、事务外访问
	和调试断点方面问题的根源。
	相比之下，传统的同步原语，包括锁和RCU，
	在同步原语和它们保护的数据之间保持清晰的分离。
\item	TM领域许多工作者的既定目标之一是
	简化大型顺序程序的并行化。
	因此，单个事务通常被期望串行执行，
	这可能在很大程度上解释了TM在多线程事务方面的问题。
\end{enumerate}

\QuickQuiz{
	鉴于像\co{spin_trylock()}这样的东西，
	声称TM引入了失败的概念怎么可能有道理？？？
}\QuickQuizAnswer{
	使用锁时，\co{spin_trylock()}是一种选择，
	相应的无失败选择是\co{spin_lock()}，
	它在通常情况下使用，即在v5.11 Linux内核中
	对\co{spin_lock()}的调用是对\co{spin_trylock()}调用的
	100多倍。
	使用TM时，唯一的无失败选择是不可撤销事务，
	它在通常情况下不被使用。
	事实上，不可撤销事务甚至并非在所有TM实现中都可用。
}\QuickQuizEnd

TM研究人员和开发者应该如何应对所有这些问题？

一种方法是关注小规模的TM，聚焦于硬件辅助可能比其他同步原语
提供显著优势的小型事务，以及有证据表明TM与锁结合的方法能够
提高生产力的小型程序~\cite{VPankratius2011TMvsLockingProductivity}。
Sun在其Rock研究CPU~\cite{DaveDice2009ASPLOSRockHTM}中
采用了小事务方法。
一些TM研究人员似乎同意这两种``小即是美''的
方法~\cite{JMStone93}，其他人对TM寄予更高期望，
而还有一些人暗示高TM抱负可能是TM最大的
敌人~\cite[Section 6]{Attiya:2010:ICT:1835698.1835699}。
尽管如此，TM完全有可能能够处理更大的问题，
而本节列出了如果TM要实现这一崇高目标
必须解决的几个问题。

当然，所有参与者都应该将此视为学习经历。
看起来TM研究人员有很多东西需要向那些
使用传统同步原语成功构建了大型软件系统的
实践者学习。

反之亦然。

\QuickQuiz{
	有什么要学的？
	为什么不直接对基于内存的数据结构使用TM，
	对那些涉及本节列出的许多愚蠢极端情况的
	罕见情况使用锁？
}\QuickQuizAnswer{
	2005年刚打来电话，它想要回它那白热化的TM营销炒作。

	在2023年，TM仍然有很多需要证明的地方，
	即使有了HTM的出现，这将在接下来的
	\cref{sec:future:Hardware Transactional Memory}中介绍。
}\QuickQuizEnd

\begin{figure}
\centering
\resizebox{3in}{!}{\includegraphics{cartoons/TM-the-vision}}
\caption{STM的愿景}
\ContributedBy{fig:future:The STM Vision}{Melissa Broussard}
\end{figure}

\begin{figure}
\centering
\resizebox{2.7in}{!}{\includegraphics{cartoons/TM-the-reality-conflict}}
\caption{STM的现实：冲突}
\ContributedBy{fig:future:The STM Reality: Conflicts}{Melissa Broussard}
\end{figure}

\begin{figure}
\centering
\resizebox{3in}{!}{\includegraphics{cartoons/TM-the-reality-nonidempotent}}
\caption{STM的现实：不可撤销操作}
\ContributedBy{fig:future:The STM Reality: Irrevocable Operations}{Melissa Broussard}
\end{figure}

\begin{figure}
\centering
\resizebox{2.7in}{!}{\includegraphics{cartoons/TM-the-reality-realtime}}
\caption{STM的现实：实时响应}
\ContributedBy{fig:future:The STM Reality: Realtime Response}{Melissa Broussard}
\end{figure}

但就目前而言，STM的当前状态最好用一系列漫画来总结。
首先，\cref{fig:future:The STM Vision}展示了STM的愿景。
一如既往，现实要更加微妙，如
\cref{fig:future:The STM Reality: Conflicts,%
fig:future:The STM Reality: Irrevocable Operations,%
fig:future:The STM Reality: Realtime Response}
所异想天开地描绘的那样。\footnote{
	最近的学术研究工作正在调查用于实时用途的基于锁的STM
	系统~\cite{JimAnderson2019STMRT,CatherineNemitz2018LockSTMrealtime}，
	尽管没有任何性能结果，并且有一些迹象表明
	实时混合STM/HTM系统必须在快速的常见情况性能和
	最坏情况前进保证之间做出
	选择~\cite{DBLP:journals/corr/AlistarhKKRS14,MartinSchoeberl2010realtimeTM}。}
不那么异想天开的STM回顾也是
可用的~\cite{JoeDuffy2010RetroTM,JoeDuffy2010RetroTM2}。

或者，有人可能会认为sequence locking构成了STM的一种
在实践中被大量使用的受限形式，如
\cref{sec:future:Case Study: Sequence Locking}中所讨论的。

无论sequence locking是否被认为是STM在实际实践中的旗手，
一些商用硬件支持HTM的受限变体，
这将在下一节中讨论。
